\section{Properties of the proportional nomination gain $\Gfun{v}{\Sagg}$}

Fix $\Sagg>0$. Recall $\Gfun{v}{\Sagg}\equiv \pOne{v}{\Sagg}/\pZero{v}{\Sagg}$ (Definition~\ref{def:Gbasic}).

The object $\Gfun{v}{\Sagg}$ isolates the \emph{primary-stage} benefit of taking action $a=1$.
Its properties are useful for (i) understanding how primary incentives vary with competitiveness,
and (ii) simplifying computation via single-index representations.


\subsection{Closed form and basic comparative statics}

\begin{lemma}[Properties of proportional nomination gain]\label{lem:Gbasic}
Fix $\Sagg > 0$. The proportional nomination gain from $a = 1$,
\[
    \Gfun{v}{\Sagg} = \frac{\pOne{v}{\Sagg}}{\pZero{v}{\Sagg}} = \frac{e^{v+1} + e\Sagg}{e^{v+1} + \Sagg},
\]
satisfies the following properties:
\begin{enumerate}
    \item $1 < \Gfun{v}{\Sagg} < e$;
    \item $\Gfun{v}{\Sagg}$ is strictly increasing in $\Sagg$ and strictly decreasing in $v$;
    \item $\Gfun{v}{\Sagg}$ depends on $(v, \Sagg)$ only through the index $v - \logg{\Sagg}$.
\end{enumerate}
\end{lemma}


Part (iii) implies that, holding fixed the opponents' profile, the primary incentive can be summarized
by a candidate's \emph{relative competitiveness} $v-\logg{\Sagg}$: higher $v$ reduces the marginal value
of taking $a=1$, while higher opponent strength (higher $\Sagg$) increases it.


\begin{proof}
Fix $\Sagg > 0$. Using $\pOneRF{v}{\Sagg}=\frac{e^{v+1}}{e^{v+1}+\Sagg}$ and
$\pZeroRF{v}{\Sagg}=\frac{e^{v}}{e^{v}+\Sagg}$, we have
\[
    \Gfun{v}{\Sagg}
    =\frac{\pOne{v}{\Sagg}}{\pZero{v}{\Sagg}}
    =\frac{e^{v+1}}{e^{v+1}+\Sagg}\cdot\frac{e^{v}+\Sagg}{e^{v}}
    =\frac{e^{v+1}+e\Sagg}{e^{v+1}+\Sagg}.
\]
\begin{enumerate}
    \item Since $\Sagg>0$ and $e^{v+1}>0$,
    \[
        \Gfun{v}{\Sagg}-1
        =\frac{e^{v+1}+e\Sagg}{e^{v+1}+\Sagg}-1
        =\frac{(e-1)\Sagg}{e^{v+1}+\Sagg}>0,
    \]
    so $\Gfun{v}{\Sagg}>1$. Moreover,
    \[
        e-\Gfun{v}{\Sagg}
        =e-\frac{e^{v+1}+e\Sagg}{e^{v+1}+\Sagg}
        =\frac{(e-1)e^{v+1}}{e^{v+1}+\Sagg}>0,
    \]
    so $\Gfun{v}{\Sagg}<e$.
    \item Differentiate the closed form:
    \[
        \frac{\partial}{\partial \Sagg}\Gfun{v}{\Sagg}
        =\frac{\partial}{\partial \Sagg}\left(\frac{e^{v+1}+e\Sagg}{e^{v+1}+\Sagg}\right)
        =\frac{(e-1)e^{v+1}}{(e^{v+1}+\Sagg)^2}>0,
    \]
    and
    \[
        \frac{\partial}{\partial v}\Gfun{v}{\Sagg}
        =\frac{\partial}{\partial v}\left(\frac{e^{v+1}+e\Sagg}{e^{v+1}+\Sagg}\right)
        =\frac{(1-e)e^{v+1}\Sagg}{(e^{v+1}+\Sagg)^2}<0.
    \]
    \item Divide numerator and denominator by $e^v$:
    \[
        \Gfun{v}{\Sagg} = \frac{e(1+\Sagg e^{-v})}{e+\Sagg e^{-v}}.
    \]
    Let $x\equiv v-\logg{\Sagg}$, so that $\Sagg e^{-v}=e^{-x}$. Then
    \[
        \Gfun{v}{\Sagg}=\frac{e(1+e^{-x})}{e+e^{-x}} \equiv \widetilde G(x),
    \]
    which yields the single-index representation.
\end{enumerate}
\end{proof}

\subsection{Derivative magnitude and competitiveness}

\begin{lemma}[Properties of $\frac{\partial}{\partial v} \Gfun{v}{\Sagg}$]\label{lem:Gderiv}
Fix $\Sagg > 0$. The derivative of $\Gfun{v}{\Sagg}$ with respect to $v$ satisfies:
\begin{enumerate}
    \item
    \[
        \frac{\partial}{\partial v}\Gfun{v}{\Sagg}
        = \frac{e^{v+1}(1-e)\Sagg}{(e^{v+1}+\Sagg)^2};
    \]
    \item the magnitude $\left|\frac{\partial}{\partial v}\Gfun{v}{\Sagg}\right|$ depends on $(v,\Sagg)$ only through the ratio $\frac{e^{v+1}}{\Sagg}$ (equivalently $\frac{e^v}{\Sagg}$ up to a constant factor);
    \item the magnitude $\left|\frac{\partial}{\partial v}\Gfun{v}{\Sagg}\right|$ is maximized when $e^{v+1} = \Sagg$.
\end{enumerate}
\end{lemma}


This lemma is useful to argue that primary incentives change most rapidly for
candidates who are neither clear frontrunners nor hopeless longshots, i.e.\ when $e^{v+1}$ and $\Sagg$
are of similar magnitude.


\begin{proof}
\begin{enumerate}
    \item This is the expression obtained in Lemma~\ref{lem:Gbasic}(ii).
    \item Using the expression from (i),
    \[
        \left|\frac{\partial}{\partial v}\Gfun{v}{\Sagg}\right|
        =(e-1)\frac{e^{v+1}\Sagg}{(e^{v+1}+\Sagg)^2}.
    \]
    Let $u\equiv \frac{e^{v+1}}{\Sagg}>0$. Then
    \[
        \left|\frac{\partial}{\partial v}\Gfun{v}{\Sagg}\right|
        =(e-1)\frac{u}{(u+1)^2},
    \]
    which depends on $(v,\Sagg)$ only through $u$.
    \item Consider $h(u)=\frac{u}{(u+1)^2}$ for $u>0$. Its derivative is
    \[
        h'(u)=\frac{1-u}{(u+1)^3},
    \]
    so $h$ is maximized at $u=1$, i.e.\ when $e^{v+1}=\Sagg$.
\end{enumerate}
\end{proof}